\documentclass[a4paper,12pt]{article}
\usepackage[utf8]{inputenc}
\usepackage[T1]{fontenc}
\usepackage[brazil]{babel}
\usepackage{geometry}
\geometry{margin=2.5cm}
\usepackage{enumitem}
\usepackage{hyperref}
\usepackage{titlesec}

\title{\textbf{Fundamentos do Design Gráfico}\\[4pt]\large Aula 01 --- Apresentação da Disciplina}
\author{Professor Pedro Sacramento}
\date{9 de junho de 2026}

\begin{document}

\maketitle

% ============================================
\section{Apresentação pessoal}

Apresentação do professor, trajetória acadêmica e profissional, áreas de atuação no design e expectativas para a disciplina.

% ============================================
\section{Sentidos}

O ser humano percebe o mundo por meio de cinco sentidos:

\begin{itemize}[nosep]
  \item Tato
  \item Olfato
  \item Visão
  \item Audição
  \item Paladar
\end{itemize}

\bigskip
\noindent O \textbf{design gráfico} concentra-se primordialmente no sentido da \textbf{visão}. É por meio da percepção visual que construímos mensagens, hierarquias visuais, identidades e experiências de comunicação.

% ============================================
\section{Exercício}

\subsection*{Objetivo}
Criar uma logo à mão e, em seguida, utilizar inteligência artificial para recriá-la com estética profissional.

\subsection*{Instruções}
\begin{enumerate}[nosep]
  \item Desenhe uma logo à mão livre (papel e caneta).
  \item Fotografe ou escaneie o desenho.
  \item Utilize o ChatGPT (ou ferramenta similar) com o seguinte \textbf{prompt}:
\end{enumerate}

\begin{quote}
\textit{``Recrie esta imagem com a estética de uma logo profissional. Dê a ela o aspecto visual de um gráfico vetorial.''}
\end{quote}

\noindent O objetivo é compreender como a IA pode ser integrada ao processo criativo, transformando um esboço conceitual em um resultado com acabamento vetorial.

% ============================================
\section{Tecnologias atuais onde conhecimentos gráficos são aplicáveis}

As competências em design gráfico estendem-se hoje a diversas tecnologias e plataformas:

\begin{itemize}[nosep]
  \item Redes sociais (criação de conteúdo visual, stories, posts, carrosséis)
  \item Impressão 3D (modelagem e texturização)
  \item Corte a laser (MDF, acrílico, papel, etc.)
  \item Sublimação (personalização de produtos)
  \item Inteligência artificial (geração de imagens, edição, automação de layout)
\end{itemize}

% ============================================
\section{Mercado}

\subsection*{Nacional}
O mercado brasileiro de design gráfico abrange desde estúdios independentes e agências de publicidade até departamentos internos em empresas de médio e grande porte. Há demanda constante por identidade visual, embalagens, materiais promocionais e conteúdo digital.

\subsection*{Internacional}
No âmbito global, plataformas de trabalho remoto e marketplaces de design (Behance, Dribbble, Fiverr, Upwork) permitem que profissionais brasileiros atuem para clientes em qualquer país, ampliando o alcance e as oportunidades de carreira.

% ============================================
\section{Formas de trabalho}

O designer gráfico pode atuar sob diferentes regimes:

\begin{itemize}[nosep]
  \item \textbf{Serviço público} --- concursos para designer, diagramador, arte-finalista em órgãos municipais, estaduais e federais.
  \item \textbf{CLT} --- contratação por empresas com carteira assinada, garantindo direitos trabalhistas e estabilidade.
  \item \textbf{PJ (Pessoa Jurídica)} --- atuação como prestador de serviços, com CNPJ próprio, emitindo notas fiscais. Oferece maior flexibilidade e remuneração potencialmente mais elevada, porém sem os benefícios da CLT.
\end{itemize}

\end{document}
